\section{Related Work}
\label{sec:related}

\subsection{Direct Affected-Version Identification}

The closest work is the large-scale benchmark study of vulnerability-affected versions~\cite{howfar2025}. It defines the task, constructs the main dataset used in this paper, and evaluates direct tools under both CVE-level and version-level metrics. Earlier work such as V-SZZ~\cite{vszz2022} explicitly connects SZZ-style vulnerability-inducing commit localization to affected-version range identification. TDSC's patch-and-developer-log based approach~\cite{tdsc2024} further highlights the importance of branch-aware reasoning and lower-bound recall when version metadata is incomplete.

Recent LLM-assisted vulnerable-version methods, including Vercation~\cite{vercation2026}, move closer to semantic patch understanding and context-aware version judgment. \toolname shares their motivation but emphasizes a stricter separation between deterministic repository planning and agent-assisted tag verdicts.

\subsection{Tracing-Based Commit Localization}

SZZ-style methods identify bug- or vulnerability-inducing commits and are often used as upstream components for affected-version inference. VCCFinder~\cite{vccfinder2015}, Lifetime~\cite{lifetime2021}, SEM-SZZ~\cite{semszz2024}, TC-SZZ~\cite{tcszz2024}, and LLM4SZZ~\cite{llm4szz2025} represent increasingly refined tracing pipelines. AgentSZZ~\cite{agentszz2026}, How and Why Agents Can Identify Bug-Introducing Commits~\cite{howwhyagents2026}, Beyond Blame~\cite{beyondblame2026}, and patch-pattern differential analysis~\cite{vicpatterns2026} show that repository exploration, pattern reasoning, and knowledge-graph or agentic inference can improve commit-level localization. However, their primary output remains a commit or commit set, whereas \toolname targets released-version state verification.

\subsection{Recurring Vulnerability and Clone-Based Detection}

Matching-based and recurring vulnerability detection systems provide useful per-version checking components. ReDeBug~\cite{redebug2012}, VUDDY~\cite{vuddy2017}, MOVERY~\cite{movery2022}, V1SCAN~\cite{v1scan2023}, FIRE~\cite{fire2024}, and VULTURE~\cite{vulture2025} detect vulnerable code clones, patch-related signatures, or reused vulnerable components. These systems demonstrate that patch-derived code evidence can be searched across versions, but they also expose the brittleness of exact or near-exact code resemblance when patches are add-only, multi-file, or semantically transformed.

\toolname treats these systems as evidence providers and baseline families rather than complete solutions to the affected-version problem. This distinction is important: a clone, hash, signature, or taint path can support a tag verdict, but the final claim still asks whether the vulnerable condition holds in a released repository state.

\subsection{Exploitability, PoC Migration, and Cross-Version Validation}

The local replication corpus also contains work on cross-version exploitability and PoC migration, including AEM~\cite{aem2023}, ARVO~\cite{arvo2024}, Diffploit~\cite{diffploit2025}, PoC migration~\cite{pocmigration2021}, library vulnerability migration~\cite{libmigration2024}, PoCGen~\cite{pocgen2025}, PwnGPT~\cite{pwngpt2025}, SemFuzz~\cite{semfuzz2017}, SyzBridge~\cite{syzbridge2024}, real-world PoC usability~\cite{pocusability2025}, and API-equivalence vulnerability detection~\cite{functionalityapi2025}. These papers are not all yet analyzed in the local \texttt{Paper/reference} corpus, so this paragraph should remain citation-placeholding until their summaries are verified \need{NEEDS CITATION VERIFICATION}.

\subsection{Agent Memory and Skill Evolution}

\toolname's long-term design includes typed memory, verifier-gated self-evolution, and reusable skills. This direction is related to agent memory surveys and systems~\cite{agentmemorysurvey2025,memps2025,memskill2026,agentsofchange2025}, as well as skill evolution and benchmarking work~\cite{autoskill2026,coevo2026,skillrl2026,skillsbench2026,trace2skill2026}. In the current paper, these references motivate future extensions rather than completed evaluation claims.
